
\chapter{Matched-Pairs Experiment}
 \label{chapter::mpe}
 
matched-pairs experiment (MPE) 是 SRE 最极端的版本：每个 stratum 中只有一个 treated unit 和一个 control unit。在这种情形下，strata 也称为 pairs。虽然这类实验是 Chapter \ref{chapter::stratification-poststratification} 中讨论的 SRE 的一个特例，但它有自己的估计和推断策略。此外，它有许多新的特点，并且与观察性研究中的 ``matching'' 策略密切相关，后者将在 Chapter \ref{chapter::matching-obs} 中讨论。因此，我们在这里用单独一章讨论 MPE。





\section{Design of the experiment and potential outcomes}
 
考虑一个包含 $2n$ 个 units 的实验。如果我们有能够预测 outcomes 的 covariates，就可以根据 covariates 的相似性来配对 units。若只有一个标量 covariate，我们可以按该 covariate 对 units 排序，然后把相邻的 units 组成 pairs。若有许多 covariates，我们可以定义 units 之间的成对距离，再基于这些距离形成 pairs。在这种情形下，pair matching 可以用 greedy algorithm 或 optimal nonbipartite matching algorithm 完成。greedy algorithm 先把距离最小的两个 units 配成一对，并将它们从 units 池中移除；再把剩余 units 中距离最小的两个配成一对，如此继续。optimal nonbipartite matching algorithm 则把 $2n$ 个 units 划分成 $n$ 个二元 pairs，使 within-pair distances 之和最小。关于 MPE 计算方面的更多细节，见 \citet{greevy2004optimal}。本章假设 pairs 已经基于 covariates 形成，然后讨论后续的设计与分析问题。


令 $(i,j)$ 表示 pair $i$ 中的 unit $j$，其中 $i=1,\ldots, n$ 且 $j=1,2$。unit $(i, j)$ 在 treatment 和 control 下分别有 potential outcomes $Y_{ij}(1)$ 和 $Y_{ij}(0)$。在每个 pair 内，我们随机分配一个 unit 接受 treatment，另一个 unit 接受 control。令
$$
Z_i = \begin{cases}
   1,& \text{if the first unit receives the treatment}, \\
   0, & \text{if the second unit receives the treatment}. 
\end{cases}  
$$ 
我们可以基于 treatment assignment mechanism 形式化地定义 MPE。


\begin{definition}
[MPE]
有
\begin{eqnarray}
 (Z_i)_{i=1}^n \iidsim \text{Bernoulli}(1/2).
 \label{eq::mp-treatment}
\end{eqnarray} 
\end{definition}

 
 
pair $i$ 内的 observed outcomes 为
$$
Y_{i1} = Z_i Y_{i1}(1) + (1-Z_i) Y_{i1}(0)
= \begin{cases}
   Y_{i1}(1),& \text{if } Z_i = 1; \\
   Y_{i1}(0), & \text{if } Z_i = 0; 
\end{cases}
$$
以及
$$
Y_{i2} = Z_i Y_{i2}(0) + (1-Z_i) Y_{i2}(1)
= \begin{cases}
   Y_{i2}(0),& \text{if } Z_i = 1; \\
   Y_{i2}(1), & \text{if } Z_i = 0 . 
\end{cases} 
$$
因此 observed data 为 $( Z_i, Y_{i1} , Y_{i2} )_{i=1}^n$。
 
 
\section{FRT}
 
与前面的讨论类似，我们总是可以用 FRT 检验 sharp null hypothesis：
$$
\HF: Y_{ij}(1) = Y_{ij}(0)  \text{ for all }  i=1,\ldots  n  \text{ and } j=1,2. 
$$
进行 FRT 时，我们需要从 \eqref{eq::mp-treatment} 模拟 $(Z_1, \ldots, Z_n)$ 的分布。
我将讨论一些基于 treated 和 control outcomes 的 within-pair differences 的经典 test statistics：
\begin{eqnarray*}
\hat{\tau}_i &=& \text{outcome under treatment} - \text{outcome under control} \text{ (within pair } i) \\
&=&  (2Z_i - 1) (Y_{i1}  - Y_{i2})\\
&=& S_i (Y_{i1}  - Y_{i2})  , 
\end{eqnarray*}
其中 $S_i = 2Z_i - 1$ 是 IID random signs，均值为 $0$、方差为 $1$，$i=1,\ldots, n$。
由于 $\hat{\tau}_i $ 为零的 pairs 不会对 randomization distribution 作出贡献，在讨论 FRT 时我们删去这些 pairs。
 
\begin{example} 
[paired $t$ statistic]\label{eg::paired-t-statistic}
within-pair differences 的平均值为
$$
\hat{\tau} = n^{-1}  \sumn \hat{\tau}_i .
$$
在 $\HF$ 下，
$$
E(\hat{\tau} ) = 0
$$
以及
$$
\var(  \hat{\tau}   ) 
= n^{-2} \sumn \var(  \hat{\tau}_i ) 
= n^{-2} \sumn \var(S_i) (Y_{i1}  - Y_{i2})^2 =  n^{-2} \sumn  \hat{\tau}_i^2 . 
$$
根据独立随机变量和的 CLT，我们有如下 Normal approximation：
$$
\frac{  \hat{\tau}  }{ \sqrt{  n^{-2} \sumn  \hat{\tau}_i^2 } } \rightarrow \N01 
$$
依分布收敛。我们可以用这个 Normal approximation 构造 asymptotic test。许多标准教科书建议在 MPE 中使用如下 paired $t$ statistic：
$$
t_\text{pair} = \frac{  \hat{\tau}  }{ \sqrt{  \{ n(n-1) \}^{-1} \sumn  (\hat{\tau}_i - \hat{\tau} ) ^2 } },
$$
当 $n$ 较大且 $\HF$ 下 $\hat{\tau} $ 较小时，它与 $\hat{\tau}$ 几乎等价。
\end{example} 

在 classical statistics 中，使用 $t_\text{pair}$ 的动机来自另一个框架。当 $\hat{\tau}_i \iidsim \textsc{N}(0, \sigma^2)$ 时，可以证明 $t_\text{pair} \sim t(n-1)$，也就是说，$t_\text{pair} $ 的 exact distribution 是自由度为 $n-1$ 的 $t$ 分布；当 $n$ 很大时，它接近 $\N01$。\ri{R} 函数 \ri{t.test} 搭配 \ri{paired=TRUE} 可以实现这个检验。当 $n$ 较大时，这些程序给出相近的结果。Example  \ref{eg::paired-t-statistic} 中的讨论在不假设数据 Normality 的情况下，为经典 paired $t$ test 提供了另一种理由。

 
 
 

 
 
  \begin{example}
 [Wilcoxon sign-rank statistic]\label{eg::mpe-sign-rank}
基于 $( | \hat{\tau}_1| , \ldots, | \hat{\tau}_n| )$ 的 ranks $(R_1, \ldots, R_n)$，我们可以定义 test statistic
$$
W = \sumn  I( \hat{\tau}_i > 0 ) R_i. 
$$
%Under $\HF$, we can drop the pairs with $ \hat{\tau}_i = 0$ since those pairs do not contribute to the randomization distribution of $W$. Ignoring those pairs with $ \hat{\tau}_i = 0$, we 
在 $\HF$ 下，$| \hat{\tau}_i|$ 是固定的，因此 $R_i$ 也是固定的，这意味着
$$
E(W) = \frac{1}{2} \sumn R_i = \frac{1}{2} \sumn i = \frac{n(n+1)}{4}
$$
以及
$$
\var(W) = \frac{1}{4} \sumn R_i^2 = \frac{1}{4} \sumn i^2 = \frac{n(n+1)(2n+1)}{24}. 
$$
独立随机变量和的 CLT 保证了如下 Normal approximation：
$$
\frac{W - n(n+1)/4}{ \sqrt{n(n+1)(2n+1) / 24    } } \rightarrow \N01
$$
依分布收敛。
我们可以用这个 Normal approximation 构造 asymptotic test。\ri{R} 函数 \ri{wilcox.test} 搭配 \ri{paired=TRUE} 可以实现 exact 和 asymptotic 两种 tests。
 \end{example}
 
 
 
 
   \begin{example}
 [Kolmogorov--Smirnov-type statistic]\label{eg::mpe-ksstatistics}
在 $\HF$ 下，absolute values $( | \hat{\tau}_1| , \ldots, | \hat{\tau}_n| )$ 是固定的，但它们的 signs 是随机的。因此 $   (  \hat{\tau}_1 , \ldots,  \hat{\tau}_n  )$ 和 $ -   (  \hat{\tau}_1 , \ldots,  \hat{\tau}_n  )$ 应该有相同的分布。
Let 
$$
\hat{F}(t)  = n^{-1} \sumn I( \hat{\tau}_i  \leq t  )
$$
为 $  (  \hat{\tau}_1 , \ldots,  \hat{\tau}_n  )$ 的 empirical distribution，并令
$$
1 - \hat{F}(-t-)  =    n^{-1} \sumn I( -  \hat{\tau}_i  \leq t  )
$$ 
为 $ -   (  \hat{\tau}_1 , \ldots,  \hat{\tau}_n  )$ 的 empirical distribution，其中 $\hat{F}(-t-)$ 是函数 $\hat{F}(\cdot)$ 在 $-t$ 处的左极限。
于是一个 Kolmogorov--Smirnov-type statistic 为
$$
D = \max_t |  \hat{F}(t)  +  \hat{F}(-t-)  - 1  |.
$$
\citet{butler1969test} 提出了这个 test statistic，并推导了它的 exact 和 asymptotic distributions。遗憾的是，标准软件包并未实现它。不过，我们可以模拟它的 exact distribution，并基于 FRT 计算 $p$-value。\footnote{
\citet{butler1969test} 在一个稍有不同的框架下提出了这个 test statistic。给定从分布 $F(y)$ 中 IID 抽取的 $    (  \hat{\tau}_1 , \ldots,  \hat{\tau}_n  )$，如果它们围绕 $0$ 对称分布，则
$$
F(t) = \pr( \hat{\tau}_i \leq t ) =  \pr( -  \hat{\tau}_i \leq t )  = 1 - \pr( \hat{\tau}_i <- t) = 1 - F(-t-).
$$
因此，$\hat{F}(t)  +  \hat{F}(-t-)  - 1 $ 衡量相对于 symmetry null hypothesis 的偏离，这也给出了 Example \ref{eg::mpe-ksstatistics} 中 $D$ 的定义动机。Kolmogorov--Smirnov-type statistic 的一个朴素定义，是像 Example \ref{eg::ks-CRE} 那样比较 treatment 和 control 下 outcomes 的 empirical distributions。使用这个定义时，我们实际上打破了 pairs。虽然它仍可用于 MPE 的 FRT，但它没有捕捉实验的 matched-pairs structure。}
 \end{example}
 
 
 
 
 
  \begin{example}
 [sign statistic]\label{eg::mpe-sign}
sign statistic 只使用 within-pair differences 的 signs：
 $$
 \Delta = \sumn I( \hat{\tau}_i > 0 ).
 $$
% Again the pairs with $\hat{\tau}_i = 0$ do not contribute to the randomization distribution of $ \Delta $, so we drop them without loss of generality. 
在 $\HF$ 下，
 $$
 I( \hat{\tau}_i > 0 ) \iidsim \textup{Bernoulli}(1/2)
 $$
因此
 $$
  \Delta \sim \textup{Binomial}(n, 1/2).
 $$
由此可得一个 exact Binomial test，它可由 \ri{R} 函数
\ri{binom.test} 搭配 \ri{p=1/2} 实现。利用 CLT，我们也可以基于 Binomial distribution 的如下 Normal approximation 进行检验：
$$
\frac{  \Delta - n/2 }{  \sqrt{n/4} } \rightarrow \N01 
$$
依分布收敛。
 \end{example}
 
 
 
 \begin{example}
 [McNemar's statistic for a binary outcome]\label{eg::mcnenar-test-binary}
如果 outcome 是 binary 的，我们可以用更紧凑的方式概括来自 MPE 的 observed data。给定一个 pair，treated outcome 可以是 $1$ 或 $0$，control outcome 也可以是 $1$ 或 $0$，于是得到 Table \ref{tb::binarymatchedpairs} 中的 two-by-two table。
 
\begin{table}
\centering
\caption{四类 pairs 的 counts}\label{tb::binarymatchedpairs}
\begin{tabular}{ccc}
\hline 
 & control outcome 1 & control outcome 0\\
 \hline 
treated outcome 1 & $m_{11}$ & $m_{10}$ \\
treated outcome 0 & $m_{01}$ & $m_{00}$ \\
\hline 
\end{tabular}
\end{table} 
 
在 $\HF$ 下，concordant pairs 的数量 $m_{11}$ 和 $m_{00}$ 是固定的，$m_{10} + m_{01}$ 也是固定的。因此唯一的随机成分是 $m_{10}$，它服从分布
$$
m_{10} \sim \textup{Binomial}( m_{10} + m_{01}, 1/2  ).
$$ 
这意味着一个基于 Binomial distribution 的 exact test。\ri{R} 函数 \ri{mcnemar.test} 给出一个基于 Binomial distribution 的 Normal approximation 的 asymptotic test：
$$
\frac{   m_{10}  - (m_{10} + m_{01})/2  }{  \sqrt{  (m_{10} + m_{01}) / 4   } } 
= \frac{ m_{10}  -  m_{01}   }{  \sqrt{ m_{10} + m_{01}    }  }
\rightarrow \N01  
$$
依分布收敛。exact FRT 和 asymptotic test 都不依赖于 $m_{11}$ 或 $m_{00}$。这些 tests 中只有 discordant pairs 的数量起作用。
 \end{example}
 
 
 
\section{Neymanian inference}
 
pair $i$ 内的 average causal effect 为
 $$
 \tau_i = \frac{1}{2}  \left\{    Y_{i1}(1) + Y_{i2}(1) -     Y_{i1}(0)  - Y_{i2}(0)     \right\},
 $$
而所有 units 的 average causal effect 为
$$
\tau = n^{-1} \sumn  \tau_i  = (2n)^{-1} \sumn \sum_{j=1}^2 \{  Y_{ij}(1) - Y_{ij}(0)   \} .
$$
直观上，$\hat{\tau}_i $ 是 $ \tau_i $ 的 unbiased estimator，因此 $\hat{\tau}  $ 是 $\tau$ 的 unbiased estimator。我们也可以计算 $\hat{\tau}  $ 的方差。由于 MPE 只是 SRE 的一个特例，我把精确公式留到 Problem \ref{para::variance-mp-design}。


然而，在 MPE 下，我们不能照搬 SRE 下的策略来估计 $\hat{\tau}  $ 的方差。outcomes 的 within-pair sample variances 并没有良好定义，因为每个 pair 内只有一个 treated unit 和一个 control unit。数据不允许我们估计 pair $i$ 内 $\hat{\tau}_i $ 的方差。


在 MPE 中是否可能估计 $\hat{\tau}$ 的方差？让我们暂时忘掉 MPE，换到经典 IID sampling 的视角。如果 $\hat{\tau}_i$ 是 IID 的，均值为 $\mu$、方差为 $\sigma^2$，那么 $\hat{\tau} = n^{-1} \sumn \hat{\tau}_i$ 的方差为 $\sigma^2 / n$。$\sigma^2$ 的一个 unbiased estimator 是 sample variance $(n-1)^{-1} \sumn ( \hat{\tau}_i -  \hat{\tau} )^2$，所以 $\var( \hat{\tau}  )$ 的一个 unbiased estimator 为
\begin{equation}\label{eq::mpe-variance-estimation}
\hat{V} = \{ n(n-1)  \}^{-1} \sumn ( \hat{\tau}_i -  \hat{\tau} )^2.
\end{equation}
上述讨论也可推广到 independent 但非 IID 的情形；见 Problem \ref{hw::inid} in Chapter \ref{chapter::basic-prob-append}。
上面的讨论看起来偏离了 MPE，因为 MPE 有完全不同的统计假设。但至少它启发了一个 variance estimator $\hat{V} $：用 $\hat{\tau}_i $ 的 between-pair variance 来估计 $\hat{\tau}  $ 的方差。当然，它是在不同假设下推导出来的。它对 MPE 也成立吗？下面的 Theorem \ref{thm::varianceest-mp} 给出了肯定答案。
 
\begin{theorem}\label{thm::varianceest-mp}
在 MPE 下，\eqref{eq::mpe-variance-estimation} 中定义的 $\hat{V}$ 是 $\hat{\tau}  $ 真实方差的 conservative estimator：
$$
E(  \hat{V}) - \var(   \hat{\tau}  ) = \{ n(n-1)  \}^{-1} \sumn  (\tau_i - \tau)^2 \geq 0  .
$$
如果 $\tau_i$ 在 pairs 之间为常数，则 $E(  \hat{V} )  = \var(\hat{\tau}) .$
\end{theorem}


Theorem \ref{thm::varianceest-mp} 说明，在 MPE 下，$\hat{V} $ 一般是 conservative variance estimator；如果 average causal effects 在 pairs 之间为常数，它就是 unbiased 的。这有些令人意外，因为 $\hat{V} $ 依赖于 $ \hat{\tau}_i $ 的 between-pair variance，而 $ \var(   \hat{\tau}  ) $ 依赖于每个 $ \hat{\tau}_i $ 的 within-pair variance。下面的证明或许能解释这个令人意外的结果。
 
 


\begin{myproof}{Theorem}{\ref{thm::varianceest-mp}}
回忆基本代数事实 $\sumn (a_i - \bar a)^2 = \sumn a_i^2 - n \bar{a}^2$，它对应于随机变量 $W$ 的恒等式 $\var(W) = E(W^2) - (EW)^2$。在下面第 2 步和第 5 步使用它，可得
\begin{eqnarray*}
n(n-1) E( \hat{V} ) &=& E\left\{   \sumn (\hat{\tau}_i - \hat{\tau})^2 \right\}  \\
&=&  E\left (    \sumn \hat{\tau}_i^2 - n \hat{\tau}^2  \right)   \\
&=& \sumn \{  \var(\hat{\tau}_i) + \tau_i^2 \} - n \{  \var(\hat{\tau})  + \tau^2 \}\\
&=& \sumn  \var(\hat{\tau}_i)  - n \var(\hat{\tau}) + \sumn  \tau_i^2 - n  \tau^2 \\
&=& n^2 \var(\hat{\tau})  - n \var(\hat{\tau})  + \sumn (\tau_i - \tau)^2.
\end{eqnarray*}
因此，
$$
E( \hat{V} )  = \var(\hat{\tau})  + \{ n(n-1)  \}^{-1}   \sumn (\tau_i - \tau)^2 \geq  \var(\hat{\tau})  .
$$
\end{myproof}




与其他实验中的讨论类似，Neymanian approach 依赖于 large-sample approximation：
$$
\frac{  \hat\tau  -  \tau }{  \sqrt{\var( \hat\tau )} } \rightarrow \N01
$$
当 $n\rightarrow \infty$ 且某些 regularity conditions 成立时，依分布收敛。由于方差被高估，Wald-type confidence interval
$$
 \hat\tau   \pm     z_{1-\alpha/2}  \sqrt{ \hat{V} }
$$
以至少 $1-\alpha$ 的概率覆盖 $\tau$。



point estimator $\hat\tau$ 和 variance estimator $\hat{V}$ 都可以方便地由 OLS 得到，如下面的 proposition 所示。

\begin{proposition}\label{prop::mpe-variance}
$\hat\tau$ 和
$\hat{V}$ 分别等于将向量 $(\hat{\tau}_1, \ldots, \hat{\tau}_n)\tran$ 只对 intercept 做 OLS fit 时，intercept 的 coefficient 和 variance estimator。
\end{proposition}

Proposition \ref{prop::mpe-variance} 的证明留作 Problem \ref{para::seols-matchedpairs}。
  
  

\section{Covariate adjustment}



虽然我们在 design stage 已经按 covariates 做了 matching，但 matching 可能并不完美 $(X_{i1}\neq X_{i2})$，而且有时我们还有 pair-matching stage 未使用的 additional covariates。在这些情形下，可以通过 adjust for the covariates 进一步提高估计效率。假设每个 unit $(i, j)$ 有 covariates $X_{ij} $。与 CRE 中的讨论类似，MPE 中有两类一般的 covariate adjustment 策略。



\subsection{FRT}

 
先从 MPE 中的 covariate-adjusted FRT 说起。与 Definition \ref{def::pseudo-y-frt} 平行，我们可以基于 outcome 对 covariates 做 model fitting 后得到的 residuals 构造 test statistics，因为在 sharp null hypothesis 下这些 residuals 是固定数。一个经典选择是把所有 observed $Y_{ij}$ 对 $X_{ij}$ 做 OLS fit，得到 residuals $\hat \varepsilon_{ij}$。然后我们可以假装 $\hat \varepsilon_{ij}$ 就是 observed outcomes 来构造 test statistics。\citet{rosenbaum2002covariance} 特别针对 MPE 提倡了这一策略。


与 Definition \ref{def::model-out-frt} 平行，我们也可以直接使用 model fitting 中的某些 coefficients 作为 test statistics。下一小节的讨论会为这一策略给出一个 test statistic 的选择。



\subsection{Regression adjustment}


现在关注 $\tau$ 的估计。我们可以像处理 outcome 一样，计算 covariates 的 within-pair differences $\hat{\tau}_{X,i} $ 及其平均值 $\hat{\tau}_X $。可以证明
$$
 E(\hat{\tau}_{X,i}   ) = 0,\quad 
 E(\hat{\tau}_X ) = 0,
 $$
以及
$$
\cov(\hat{\tau}_X   ) =n^{-2} \sumn \hat{\tau}_{X,i}   \hat{\tau}_{X,i}  \tran .
%= n^{-2} \sumn |  \hat{\tau}_{X,i}  |  |\hat{\tau}_{X,i}  | \tran  .
$$
在一个 realized MPE 中，除非所有 $\hat{\tau}_{X,i}  $ 都为零，否则 $\cov(\hat{\tau}_X   ) $ 不为零。如果 $(Z_1,\ldots, Z_n)$ 抽得不走运，$\hat{\tau}_X $ 可能明显偏离零。与 Chapter \ref{sec::lin=adjustingforimbalance} 中 CRE 下的讨论类似，adjusting for the imbalance of the covariate means 很可能提高估计效率。

与 Section \ref{sec::lin=adjustingforimbalance} 中的讨论类似，我们可以考虑一类由 $\gamma$ 标记的 estimators：
$$
\hat{\tau}(\gamma) = \hat{\tau} - \gamma\tran \hat{\tau}_X 
$$
对于任意固定的 $\gamma$，它的均值为 $0$。我们希望选择 $\gamma$ 来最小化 $\hat{\tau}(\gamma)$ 的方差。它的方差是 $\gamma$ 的二次函数：
\begin{eqnarray*}
\var\{ \hat{\tau}(\gamma) \} &=&   \var(\hat{\tau} - \gamma\tran \hat{\tau}_X ) \\
&=&  \var(\hat{\tau}  ) + \gamma\tran \cov(\hat{\tau}_X   ) \gamma - 2\gamma\tran \cov( \hat{\tau}_X ,  \hat{\tau}  ),
\end{eqnarray*}
其最小值在如下位置取得：
$$
\tilde{\gamma} = \cov(\hat{\tau}_X   )^{-1}  \cov( \hat{\tau}_X ,  \hat{\tau}  ).
$$
上面已经得到了 $\cov(\hat{\tau}_X   ) $ 的公式，它也可以写作
$$
\cov(\hat{\tau}_X   ) = n^{-2} \sumn |  \hat{\tau}_{X,i}  |  |\hat{\tau}_{X,i}  | \tran,
$$
其中 $| \cdot | $ 表示向量逐分量取绝对值。因此 $\cov(\hat{\tau}_X   ) $ 是固定的，并且可由 observed data 得知。然而，$\cov( \hat{\tau}_X ,  \hat{\tau}  )$ 依赖于未知的 potential outcomes。幸运的是，我们可以为它得到一个 unbiased estimator，如下面的 Theorem \ref{thm::covariance-mp} 所示。

\begin{theorem}\label{thm::covariance-mp}
$\cov( \hat{\tau}_X ,  \hat{\tau}  )$ 的一个 unbiased estimator 为
$$
\hat{\theta} = \{n(n-1)  \}^{-1} \sumn ( \hat{\tau}_{X,i}  -  \hat{\tau}_{X} ) (   \hat{\tau}_i -  \hat{\tau}   ).
$$
\end{theorem}

Theorem \ref{thm::covariance-mp} 的证明类似于 Theorem \ref{thm::varianceest-mp} 的证明。我把它留作 Problem \ref{para::covariance-est-proof}。


因此，我们可以用下面的式子估计 optimal coefficient $\tilde{\gamma}$：
\begin{eqnarray*}
\hat{\gamma} &=&  \left( n^{-2} \sumn \hat{\tau}_{X,i}   \hat{\tau}_{X,i}  \tran  \right)^{-1}
\left\{ \{n(n-1)  \}^{-1} \sumn (\hat{\tau}_{X,i}  -  \hat{\tau}_X )(   \hat{\tau}_i -  \hat{\tau}   ) \right\} \\
&\approx &  \left(  \sumn(  \hat{\tau}_{X,i}   - \hat{\tau}_X ) ( \hat{\tau}_{X,i}  -  \hat{\tau}_X  )\tran  \right)^{-1}
  \sumn (\hat{\tau}_{X,i}  -  \hat{\tau}_X ) (   \hat{\tau}_i -  \hat{\tau}   )  ,
\end{eqnarray*}
它近似等于在带 intercept 的 OLS 中，把 $\hat{\tau}_i $ 对 $\hat{\tau}_{X,i} $ 回归时 $\hat{\tau}_{X,i}$ 的 coefficient。最终 estimator 为
$$
\hat{\tau}_{\text{adj}} = 
\hat{\tau}( \hat{\gamma} ) = \hat{\tau} - \hat{\gamma} \tran \hat{\tau}_X  ,
$$
由 OLS 的性质，它近似等于在带 intercept 的 OLS 中，把 $\hat{\tau}_i $ 对 $\hat{\tau}_{X,i} $ 回归时的 intercept。


于是 $\hat{\tau}_{\text{adj}} $ 的一个 conservative variance estimator 为
\begin{eqnarray*}
\hat{V} _{\text{adj}} 
&=& 
\hat{V} + \hat \gamma\tran \cov(\hat{\tau}_X   ) \hat \gamma - 2 \hat \gamma\tran \hat \theta  \\
&=& \hat{V}  - \hat \theta \tran \cov(\hat{\tau}_X   ) ^{-1}    \hat \theta .
\end{eqnarray*}


一个微妙的技术问题是，$\hat\tau(\hat\gamma)$ 是否与 $\hat\tau(\tilde\gamma)$ 具有相同的 optimality。我们在讨论 \citet{lin2013} 的 estimator 时遇到过类似问题。
在大样本下，可以证明 $\hat\tau(\hat\gamma) - \hat\tau(\tilde\gamma) =  - (\hat\gamma - \tilde\gamma)\tran \hat{\tau}_X $ 是高阶项，因为它是两个 ``small'' terms $\hat\gamma - \tilde\gamma$ 和 $\hat{\tau}_X$ 的乘积。


此外，\citet{fogarty2018regression} 讨论了上述 covariate-adjusted procedure 的渐近等价 regression formulation，并给出了相关 CLT 的严格证明。我在下面概括这个 regression formulation，但不列出 regularity conditions。


\begin{proposition}\label{prop::colin2018regression}
在 MPE 下，covariate-adjusted estimator $\hat{\tau}_{\textup{adj}} $ 及其 associated variance estimator $\hat{V} _{\textup{adj}}$ 可以方便地近似为如下 OLS fit 中的 intercept 及其 associated variance estimator：把 $\hat{\tau}_i $ 的向量对 $1$ 向量和 $\hat{\tau}_{X,i} $ 矩阵做 OLS fit。
\end{proposition}


Proposition \ref{prop::colin2018regression} 的证明留作 Problem \ref{para::seols-matchedpairs}。有意思的是，Proposition \ref{prop::mpe-variance} 和 \ref{prop::colin2018regression} 都不需要对 variance estimator 做 EHW correction。更多技术细节见 \citet{fogarty2018regression}。由于我们把 MPE 的数据化约为 within-pair differences，因此没有必要 center covariates；这不同于 \citet{lin2013} 的 estimator 在 CRE 下的实现方式。

 


\section{Examples}

\subsection{Darwin's data comparing cross-fertilizing and self-fertilizing on the height of corns}
\label{sec::frt-darwin-mpe}

这是来自 \citet{Fisher:1935} 的一个经典例子。它包含 15 对玉米，分别进行 cross-fertilizing 或 self-fertilizing，height 是 outcome。\ri{R} package \ri{HistData} 提供原始数据，其中 \ri{cross} 和 \ri{self} 分别是在 cross-fertilizing 和 self-fertilizing 下的 heights，\ri{diff} 表示二者的差。

\begin{lstlisting}
> library("HistData")
> ZeaMays
   pair pot  cross   self   diff
1     1   1 23.500 17.375  6.125
2     2   1 12.000 20.375 -8.375
3     3   1 21.000 20.000  1.000
4     4   2 22.000 20.000  2.000
5     5   2 19.125 18.375  0.750
6     6   2 21.500 18.625  2.875
7     7   3 22.125 18.625  3.500
8     8   3 20.375 15.250  5.125
9     9   3 18.250 16.500  1.750
10   10   3 21.625 18.000  3.625
11   11   3 23.250 16.250  7.000
12   12   4 21.000 18.000  3.000
13   13   4 22.125 12.750  9.375
14   14   4 23.000 15.500  7.500
15   15   4 12.000 18.000 -6.000
\end{lstlisting}
 


总共来说，这个 MPE 有 $2^{15} = 32768$ 种 possible treatment assignment，这个数量在 \ri{R} 中是可处理的。下面的函数可以枚举 MPE 的所有 possible treatment assignments：

\begin{lstlisting}
MP_enumerate = function(i, n.pairs) 
{
 if(i > 2^n.pairs)  print("i is too large.")
 a = 2^((n.pairs-1):0)
 b = 2*a
 2*sapply(i-1, 
          function(x) 
            as.integer((x %% b)>=a)) - 1
}
\end{lstlisting}



于是我们枚举所有 treatment assignments，并计算相应的 $\hat{\tau}$ 以及 one-sided exact $p$-value。

\begin{lstlisting}
> difference = ZeaMays$diff 
> n.pairs    = length(difference)
> abs.diff   = abs(difference)
> t.obs      = mean(difference)
> t.ran      = sapply(1:2^15, 
+                     function(x){ 
+                       sum(MP_enumerate(x, 15)*abs.diff) 
+                       })/n.pairs
> pvalue     = mean(t.ran>=t.obs)
> pvalue
[1] 0.02633667
\end{lstlisting}

Figure \ref{fig::frt-darwin-mpe} 展示了 $\HF$ 下 $\hat{\tau}$ 的 exact randomization distribution。

\begin{figure}[th]
\centering
\includegraphics[width = \textwidth]{figures/FRT_Darwin.pdf}
\caption{使用 Darwin's data 得到的 $\hat{\tau}$ 的 exact randomization distribution}\label{fig::frt-darwin-mpe}
\end{figure}



\subsection{Children's television workshop experiment data}
\label{sec::neyman-television-mpe}



我还重新分析 \citet{ball1973reading} 中 Children's Television Workshop experiment 的一个数据子集；\citet[][Chapter 10]{imbens2015causal} 也分析过该数据。它包含 8 对 classes，每一对中的一个 class 被随机分配在标准 reading-class period 中观看 {\it The Electric Company} 节目。数据包含 pre-test score（covariate）和 post-test score（outcome）。
下面的表概括了 within-pair covariates 和 outcomes 及其 differences：
\begin{lstlisting}
> dataxy = c(12.9, 12.0, 54.6, 60.6,
+            15.1, 12.3, 56.5, 55.5,
+            16.8, 17.2, 75.2, 84.8,
+            15.8, 18.9, 75.6, 101.9,
+            13.9, 15.3, 55.3, 70.6,
+            14.5, 16.6, 59.3, 78.4,
+            17.0, 16.0, 87.0, 84.2,
+            15.8, 20.1, 73.7, 108.6)
> dataxy = matrix(dataxy, 8, 4,  byrow = TRUE)           
> diffx = dataxy[, 2] - dataxy[, 1]
> diffy = dataxy[, 4] - dataxy[, 3]
> dataxy = cbind(dataxy, diffx, diffy)
> rownames(dataxy) = 1:8
> colnames(dataxy) = c("x.control", "x.treatment", 
+                      "y.control", "y.treatment",
+                      "diffx", "diffy")
> dataxy = as.data.frame(dataxy)
> dataxy
  x.control x.treatment y.control y.treatment diffx diffy
1      12.9        12.0      54.6        60.6  -0.9   6.0
2      15.1        12.3      56.5        55.5  -2.8  -1.0
3      16.8        17.2      75.2        84.8   0.4   9.6
4      15.8        18.9      75.6       101.9   3.1  26.3
5      13.9        15.3      55.3        70.6   1.4  15.3
6      14.5        16.6      59.3        78.4   2.1  19.1
7      17.0        16.0      87.0        84.2  -1.0  -2.8
8      15.8        20.1      73.7       108.6   4.3  34.9
\end{lstlisting}


下面的 \ri{R} 代码计算 $\hat\tau$ 和 $\hat V$。
\begin{lstlisting}
> n      = dim(dataxy)[1]
> tauhat = mean(dataxy[, "diffy"])
> vhat   = var(dataxy[, "diffy"])/n
> tauhat
[1] 13.425
> sqrt(vhat)
[1] 4.636337
\end{lstlisting}
根据 Propositions \ref{prop::mpe-variance} 和 \ref{prop::colin2018regression}，我们可以使用 OLS 得到 point estimators 和 standard errors。不调整 covariates 时，有
\begin{lstlisting}
> unadj = summary(lm(diffy ~ 1, data = dataxy))$coef
> round(unadj, 3) 
            Estimate Std. Error t value Pr(>|t|)
(Intercept)   13.425      4.636   2.896    0.023
\end{lstlisting}
调整 covariates 后，有
\begin{lstlisting} 
> adj = summary(lm(diffy ~ diffx, data = dataxy))$coef
> round(adj, 3)
            Estimate Std. Error t value Pr(>|t|)
(Intercept)    8.994      1.410   6.381    0.001
diffx          5.371      0.599   8.964    0.000
\end{lstlisting}


上述结果假设 $n$ 很大；如果我们相信 large-$n$ approximation，那么 $p$-values 是有依据的。然而，$n=8$ 并不大。总共有 $2^8 = 256$ 种 possible treatment assignments，因此最小可能的 $p$-value 是 $1/256 = 0.0039$，这远大于 covariate-adjusted estimator 的 Normal approximation 给出的 $p$-value。在这个例子中，更合理的做法是使用带 studentized statistic 的 FRT 来计算 exact $p$-values；这里 studentized statistic 是 \ri{lm} 函数给出的 \ri{t value}\footnote{理由见 Chapter \ref{chapter::unification-fisher-neyman}。}。下面的 \ri{R} 代码基于两个 $t$-statistics 计算 $p$-values。
\begin{lstlisting} 
> t.ran = sapply(1:2^8, function(x){ 
+   z.mpe = MP_enumerate(x, 8)
+   diffy.mpe = diffy*z.mpe
+   diffx.mpe = diffx*z.mpe
+   
+   c(summary(lm(diffy.mpe ~ 1))$coef[1, 3],
+     summary(lm(diffy.mpe ~ diffx.mpe))$coef[1, 3]) 
+ })
> p.unadj = mean(abs(t.ran[1, ]) >= abs(unadj[1, 3]))
> p.unadj
[1] 0.03125
> p.adj = mean(abs(t.ran[2, ]) >= abs(adj[1, 3]))
> p.adj
[1] 0.0078125
\end{lstlisting}


Figure \ref{fig::frt-television-mpe} 展示了两个 studentized statistics 的 exact distributions 以及 two-sided $p$-values。该图突出了一个事实：test statistics 的 randomization distributions 是离散的，最多只取 256 个可能值。Normal approximations 很可能不准确，尤其是在尾部。我们应当报告基于 FRT 的 $p$-values。


\begin{figure}[th]
\centering
\includegraphics[width = \textwidth]{figures/frt_television.pdf}
\caption{Section \ref{sec::neyman-television-mpe} 中 studentized statistics 的 exact randomization distributions}\label{fig::frt-television-mpe}
\end{figure}



\section{Comparing the MPE and CRE}
\label{sec::discussion}

\citet{imai2008variance} 比较了 MPE 和 CRE。启发式地说，结论是：如果 matching 做得好，并且 covariates 能预测 outcome，那么 MPE 会给出更精确的 estimators。然而，在 design stage 没有 outcome data 时，很难判断这一点是否成立。在 FRT 中，如果 covariates 能预测 outcome，那么与 CRE 相比，MPE 通常给出更有 power 的 tests。\citet{greevy2004optimal} 用基于 Wilcoxon sign-rank statistic 的 simulation 说明了这一点。不过，在 finite samples 下，这可能是一个微妙问题。考虑一个包含 $2n$ 个 units 的实验，其中 $n$ 个 units 接受 treatment，$n$ 个 units 接受 control。如果我们在 level $0.05$ 下检验 sharp null hypothesis，那么在 MPE 中至少需要 $2\times 5=10$ 个 units，因为最小 $p$-value 是 $1/2^5 = 1/32 < 0.05$，而 $1/2^4 = 1/16 > 0.05$；但在 CRE 中至少只需要 $2\times 4 = 8$ 个 units，因为最小 $p$-value 是 $1/\binom{8}{4} = 1/70 < 0.05$，而 $1/\binom{6}{3} = 1/20 = 0.05$。因此，当有 $8$ 个 units 时，在 MPE 中不可能拒绝 sharp null hypothesis，但在 CRE 中可能做到。即便 covariates 是 outcome 的完美预测量，基于 FRT 来看，MPE 也并不优于 CRE。



\section{Extension to the general matched experiment}\label{sec::general-matched-experiment}

把 MPE 推广到 control units 数量可变的 general matched experiment 是直接的。假设我们有 $n$ 个 matched sets，用 $i=1, \ldots, n$ 标记。对于 matched set $i$，有 $1+M_i$ 个 units。$M_i$ 可以变化。experimental units 的总数为 $N = n + \sum_{i=1}^n M_i$。
令 $ij$ 表示 matched set $i$ 内的 unit $j$（$i=1,\ldots, n$ 且 $j=1,\ldots, M_i+1$）。unit $ij$ 在 treatment 和 control 下分别有 potential outcomes $Y_{ij}(1)$ 和 $Y_{ij}(0)$。


在 matched set $i$ $(i=1, \ldots, n)$ 内，实验者随机选择且只选择一个 unit 接受 treatment，其余 $M_i$ 个 units 接受 control。
这个 general matched experiment 也是 SRE 的一个特例：它有 $n$ 个 strata，第 $i$ 个 stratum 的大小为 $1+M_i$ $(i=1, \ldots, n)$。
令 $Z_{ij}$ 为 unit $ij$ 的 treatment indicator，它揭示其中一个 potential outcome：
$$
Y_{ij} = Z_{ij} Y_{ij}(1) + (1-Z_{ij}) Y_{ij}(0) . 
$$ 

matched set $i$ 内的 average causal effect 等于
$$
\tau_i = (M_i + 1)^{-1} \sum_{j=1}^{1+M_i} \{ Y_{ij}(1)-Y_{ij}(0) \}. 
$$
由于 general matched experiment 是 SRE，$\tau_i $ 的一个 unbiased estimator 为
$$
\hat \tau_i = \sum_{j=1}^{M_i + 1} Z_{ij} Y_{ij} - M_i^{-1} \sum_{j=1}^{M_i+1} (1-Z_{ij}) Y_{ij}
$$
它就是 matched set $i$ 内 outcomes 的 difference in means。

下面讨论 general matched experiment 下的 statistical inference。


\subsection{FRT}

和往常一样，我们总是可以用 FRT 检验 sharp null hypothesis：
$$
\HF: Y_{ij}(1) = Y_{ij}(0)  \text{ for all }  i=1,\ldots, n; j=1,\ldots, M_i+1. 
$$
因为 general matched experiment 是具有许多小 strata 的 SRE 的一个特例，我们可以使用 Examples \ref{eg::hodgeslehmann}, \ref{eg::ks-stratified}, \ref{eg::mpe-sign-rank}, \ref{eg::mpe-ksstatistics}, \ref{eg::mpe-sign} 中定义的 test statistics，也可以使用下面两个小节中的 estimators 及其对应的 $t$-statistics。



\subsection{Estimating the average of the within-strata effects}


我们首先关注 within-strata effects 的平均值：
$$
\tau = n^{-1} \sumn \tau_i .
$$
它有一个 unbiased estimator：
$$
\hat\tau = n^{-1} \sumn \hat \tau_i . 
$$
有意思的是，可以证明 Theorem \ref{thm::varianceest-mp} 对 general matched experiment 也成立，MPE 的其他结果也同样成立。特别地，我们可以把 $\hat \tau_i $ 对 intercept 做 OLS fit，从而得到 $\tau$ 的 point 和 variance estimators。有 covariates 时，可以把 $\hat \tau_i $ 对 intercept 和 $\hat\tau_{X,i}$ 做 OLS fit，其中
$$
\hat \tau_{X,i} = \sum_{j=1}^{M_i + 1} Z_{ij} X_{ij} - M_i^{-1} \sum_{j=1}^{M_i+1} (1-Z_{ij}) X_{ij}
$$
是 matched set $i$ 内 covariates 的相应 difference in means。



\subsection{A more general causal estimand}


重要的是，上面的 $\tau$ 是 $\tau_i$ 的平均值；当 $M_i$ 变化时，它不等于实验中 $N$ 个 units 的 average causal effect。average causal effect 等于
$$
\tau' = N^{-1} \sumn \sum_{j=1}^{1+M_i} \{ Y_{ij}(1)-Y_{ij}(0) \} = \sumn \frac{1+M_i}{N} \tau_i.
$$ 
为了统一讨论，我考虑 weighted causal effect：
$$
\tau_w =  \sumn w_i \tau_i
$$
with $\sumn w_i = 1$, 
它包含 $\tau$ 作为特例（$w_i = n^{-1}$），也包含 $\tau'$ 作为特例（$i=1,\ldots, n$ 时 $w_i = (1+M_i)/N $）。容易得到一个 unbiased estimator：
$$
\hat\tau_w =  \sumn w_i \hat\tau_i,
$$
并计算其方差：
$$
\var(\hat\tau_w ) = \sumn w_i^2 \var (\hat\tau_i ). 
$$
然而，估计 $\hat\tau_w $ 的方差相当微妙，因为 $\hat\tau_i$ 是没有 replicates 的 independent random variables。这是 theoretical statistics 中的一个著名问题，\citet{hartley1969variance} 和 \citet{rao1970estimation} 研究过它。\citet{fogarty2018mitigating} 也讨论了这个问题，但没有提及这些早期工作。我将直接给出 variance estimator 的最终形式，而不详细说明其动机：
$$
\hat{V}_w = \sumn c_i (\hat\tau_i -\hat\tau_w )  ^2
$$
其中
$$
c_i = \frac{    \frac{w_i^2}{1-2w_i}   }{   1 + \sumn  \frac{w_i^2}{1-2w_i} }.
$$
作为 sanity check，在 $M_i=1$ 且 $w_i = n^{-1}$ 的 MPE 中，$c_i $ 化为 $ \{  n(n-1) \}^{-1}$。为简单起见，我们关注所有 $i$ 都满足 $w_i < 1/2$ 的情形，也就是说，没有任何 matched set 包含超过一半的总权重。下面的 theorem 推广了 Theorem \ref{thm::varianceest-mp}。

\begin{theorem}
\label{thm::hadamdard-matched-experiment}
在 $M_i$ 可变的 general matched experiment 下，如果所有 $i$ 都有 $w_i < 1/2$，则
 $$
 E(\hat{V}_w ) - \var( \hat\tau_w  ) = \sumn c_i (\tau_i - \tau_w )^2 \geq \var( \hat\tau_w  )  \geq 0
 $$
当 $\tau_i $ 为常数时等号成立。
\end{theorem}

虽然 $\hat{V}_w $ 的理论动机相当复杂，但直接验证 Theorem \ref{thm::hadamdard-matched-experiment} 并不太困难。我把证明留到 Problem \ref{para::variance-general-matched}。
 



\section{Homework Problems}

\homeworksolutionref{07}
\paragraph{The true variance of $\hat{\tau}$ in the MPE}
\label{para::variance-mp-design}

用 potential outcomes 的 finite-population variances 表示 $\var(\hat{\tau})$。

\paragraph{A covariance estimator}\label{para::covariance-est-proof}
证明 Theorem \ref{thm::covariance-mp}。

\paragraph{Variance estimators via OLS}\label{para::seols-matchedpairs}
证明 Propositions \ref{prop::mpe-variance} 和 \ref{prop::colin2018regression}。


\paragraph{Point and variance estimator with binary outcome}\label{para::binary-matchedpairs-neyman}


本题把 Example \ref{eg::mcnenar-test-binary} 推广到 Neymanian inference。

用 Table \ref{tb::binarymatchedpairs} 中的 counts 表示 $\hat\tau$ 和 $\hat V$。




\paragraph{Minimum sample size for the FRT}
\label{prob::minimum-n-0.001}

推广 Section \ref{sec::discussion} 中的讨论。考虑一个包含 $2n$ 个 units 的实验，其中 $n$ 个 units 接受 treatment，$n$ 个 units 接受 control，并在 level $0.001$ 下检验 sharp null hypothesis。对于 MPE，要使最小 $p$-value 不超过 $0.001$，$n$ 的最小值是多少？CRE 中对应的 $n$ 的最小值又是多少？
 

\paragraph{Re-analyzing Darwin's data}

 
Chapter \ref{sec::frt-darwin-mpe} 使用基于 test statistic $\hat{\tau}$ 的 FRT 分析了 Darwin's data。

使用带 Wilcoxon sign-rank statistic 的 FRT 重新分析该数据集。

基于 Neymanian inference 重新分析该数据集：报告 unbiased point estimator、conservative variance estimator 和 95\% confidence interval。

\paragraph{Re-analyzing children's television workshop experiment data}

Chapter \ref{sec::neyman-television-mpe} 基于 Neymanian inference 分析了该数据。

使用带不同 test statistics 的 FRT 重新分析该数据集。

使用带 covariate adjustment 的 FRT 重新分析该数据集。




\paragraph{Re-analyzing \citet{angrist2009effects}'s data}\label{para::AL-mpe-neyman}

\citet{angrist2009effects} 的原始分析相当复杂。本题只关注原论文的 Table A1，并把 schools 视为 experimental units。\citet{angrist2009effects} 实质上是在 schools 上实施了一个 MPE。删去 pair 6 以及所有存在 noncompliance 的 pairs 后，得到 14 个 complete pairs，数据如下，也见 \ri{AL2009.csv}：
\begin{lstlisting}
   pair z  pr99  pr00  pr01  pr02
1     1 0 0.046 0.000 0.091 0.185
2     1 1 0.036 0.051 0.000 0.047
3     2 0 0.054 0.094 0.184 0.034
4     2 1 0.050 0.108 0.110 0.095
5     3 0 0.114 0.000 0.056 0.075
6     3 1 0.098 0.054 0.030 0.068
7     4 0 0.148 0.162 0.082 0.075
8     4 1 0.134 0.390 0.339 0.458
9     5 0 0.152 0.105 0.083 0.129
10    5 1 0.145 0.077 0.579 0.167
11    6 0 0.188 0.214 0.375 0.545
12    6 1 0.179 0.165 0.483 0.444
13    7 0 0.193 0.771 0.328 0.583
14    7 1 0.189 0.186 0.168 0.368
15    8 0 0.197 0.350 0.000 0.383
16    8 1 0.200 0.071 0.667 0.429
17    9 0 0.213 0.176 0.164 0.172
18    9 1 0.209 0.165 0.092 0.151
19   10 0 0.211 0.667 0.250 0.617
20   10 1 0.219 0.250 0.500 0.350
21   11 0 0.219 0.153 0.185 0.219
22   11 1 0.224 0.363 0.372 0.342
23   12 0 0.255 0.226 0.213 0.327
24   12 1 0.257 0.098 0.107 0.095
25   13 0 0.261 0.071 0.000    NA
26   13 1 0.263 0.441 0.448 0.435
27   14 0 0.286 0.161 0.126 0.181
28   14 1 0.285 0.389 0.353 0.309
\end{lstlisting}
outcomes 是 2001 年和 2002 年的 Bagrut passing rates，1999 年和 2000 年的 Bagrut passing rates 作为 pretreatment covariates。基于 Neymanian inference，在使用和不使用 covariates 的情况下重新分析该数据。特别地，你如何处理 pair 25 中缺失的 outcome？
% (1) using the FRT with and without covariate adjustment and (2) based on the Neymanian inference with and without covariates. In particular, how do you deal with the missing outcome? 
 
 
\paragraph{Variance estimation in the general matched experiment} 
 \label{para::variance-general-matched}

本题补充 Section  \ref{sec::general-matched-experiment} 的更多细节。
 
 
首先，为 general matched experiment 证明 Theorem  \ref{thm::varianceest-mp}。
其次，证明 Theorem \ref{thm::hadamdard-matched-experiment}。
 
Remark: 对于第二部分，可以先验证 $\hat\tau_i -\hat\tau_w$ 的均值为 $\tau_i - \tau_w$，方差为
 $$
 \var(\hat\tau_i -\hat\tau_w) = \var(\hat\tau_w ) + (1-2w_i) \var(\hat\tau_i). 
 $$


\paragraph{Recommended readings}

\citet{greevy2004optimal} 提供了一种基于 covariates 形成 matched pairs 的算法。
\citet{imai2008variance} 讨论了不使用 covariates 时 average causal effect 的估计，
而
\citet{fogarty2018regression} 讨论了 MPEs 中的 covariate adjustment。
